\subsection{Conclusions}

The quantitative analysis of \ac{HLT} and phase composition in hydrated \ac{C3S}, \ac{C2S}, and their mixtures demonstrated the potential of image processing beyond mere segmentation.
The method can be utilised to determine the \ac{HLT} development in \ac{C3S}, \ac{C2S}, and \ac{C3S}:\ac{C2S} mixtures up to \qty{28}{\day} of hydration.
The study confirmed the faster hydration of \ac{C3S} and revealed a slight retardation of \ac{C3S} hydration in mixtures with high \ac{C2S} content.
Furthermore, the formation of \textsc{Hadley} grains was observed in mixtures with higher \ac{C2S} content.
Another observation was that clinker particles enclosed by dense \ac{CH} show inhibited hydration, as the lack of adjacent pore space prevents the dissolution and the resulting growth of \ac{CSH}.
Finally, the assessment of statistical representativeness indicated that the required analysis area is dictated by microstructural heterogeneity, ranging from \qty{0.3}{\milli\meter\squared} for homogeneous samples to over \qty{1.4}{\milli\meter\squared} for those with significant irregularities such as air voids or inconsistent \ac{CH} growth.

The image evaluation method developed for \ac{HLT} measurements was also adapted to quantify inter-particle distances in fresh cement paste from cryo-\ac{FIB}-\ac{SEM} images.
Applied to \ac{C3S} paste, \ac{LC3} with and without \ac{PUS}, and to published \ac{OPC}+\ac{PCE} data, the histograms show fewer large inter-particle distances after dispersion.
These large distances are water-filled spaces between agglomerates.
Their reduction is the microstructural signature resolved by this measurement.
For \ac{PUS} applied to \ac{LC3}, the decrease in the mean distance still lies within the uncertainty of the measurement.
The broader distribution of particle distances in untreated specimens, compared to sonicated or superplasticised ones, is reflected in the histograms and supports the qualitative impression from the images.
The main limitation remains the small imaged area of cryo-\ac{FIB} sections, which restricts statistical representativeness and particle counts. 
The approach is nevertheless suitable to objectify dispersion effects in fresh paste and to complement rheological or other bulk measurements.

Finally, the correlation of the \ac{HLT} measurements with \ac{EBSD} data investigated whether \ac{CSH} grows preferentially on specific crystallographic orientations of the \ac{C3S} surface.
Projecting the measured \ac{HLT}s onto a unit sphere representing the \ac{C3S} crystal orientation and analysing the distribution revealed a fairly homogeneous \ac{HLT} distribution.
Despite some local variations associated with low measurement density, the current results indicate that there is no preferred crystallographic direction for the growth of \ac{CSH} around \ac{C3S} particles.
However, to provide a more definitive conclusion, further investigations with higher \ac{EBSD} resolutions, larger analysis areas, and an improved segmentation distinguishing between \ac{CSH} and \ac{CH} are recommended.

%Machine learning-based pixel classification proved superior to thresholding, providing accuracy comparable to manual segmentation.